\documentclass[a4paper, twocolumn, 11pt, twoside]{article}

\usepackage[portuges]{babel}
\usepackage{linguamatica}
\usepackage[T1]{fontenc}
\usepackage{amsmath,amssymb}
\usepackage[safe]{tipa}

\bibliographystyle{sp_por}

\submitted{17 de \NOV{} de 2026}
\accepted{---}

\title{LipSpell em Português Brasileiro: Correção Ortográfica Fonética com Traços Distintivos, Regras G2P e Análise de Erros}
\titleEN{LipSpell in Brazilian Portuguese: Phonetic Spelling Correction with Distinctive Features, G2P Rules, and Error Analysis}

\author{
  Leonardo Araujo
  \instituto{Universidade Federal de São João del-Rei (UFSJ)}
  \email{leolca@ufsj.edu.br}
}

\begin{document}
\maketitle

\begin{resumo}
Corretores ortográficos tradicionais baseiam-se em distância de edição em nível de grafema e proximidade de teclado --- falhando em erros cognitivos e fonéticos --- ou em codificações fonéticas específicas para o inglês, como Soundex e Metaphone. Neste trabalho, detalhamos, adaptamos e avaliamos em profundidade para o português brasileiro o \textbf{LipSpell}, um arcabouço de pontuação fonética parametrizável por idioma proposto originalmente em trabalho paralelo, integrando: (1) uma matriz pré-computada $|P|\times|P|$ de distâncias entre fones derivada de traços distintivos universais do Alfabeto Fonético Internacional (IPA), (2) geometria física do teclado ABNT2 com tolerância a diacríticos via decomposição canônica Unicode NFD (\textit{Normalization Form D}) e (3) probabilidades a priori de frequência lexical, acelerado por indexação em dois estágios (remoções simétricas \textit{SymDelete} e código fonético de sete classes). Na base \textit{Wikipédia: Lista de erros comuns} ($n=1.149$ pares e fatia $n=500$), o LipSpell supera Norvig, Hunspell e Aspell ($73{,}20\%$ na transferência direta e $76{,}20\%$ com tolerância diacrítica NFD no teclado vs.\ $66{,}40\%$/$64{,}20\%$/$64{,}00\%$ em $n=500$; $72{,}24\%$ e $74{,}93\%$ vs.\ $66{,}23\%$/$63{,}19\%$/$60{,}14\%$ em $n=1.149$), com ganhos expressivos em erros apenas de diacrítico ($+11{,}62$ pontos percentuais sobre Norvig) e múltiplas edições ($+15{,}34$ pp). Comparamos ainda o conversor G2P baseado em regras do português (\texttt{PortugueseG2P}) contra três dicionários de pronúncia (ABG, \texttt{ipa-dict}/Wiktionary e WikiPron) no pipeline completo: com o refinamento de vogais graves (\texttt{è/ì/ò/ù}) e dessonorização de \texttt{-z} final, as regras puras igualam o ABG em Top-1 nas fatias completas ($76{,}20\%$ em $n=500$; $74{,}93\%$ em $n=1.149$), alcançam o maior Top-5 nessas fatias ($85{,}20\%$ e $83{,}90\%$) e superam os três dicionários em divisões limpas de erros \textit{non-word} ($87{,}17\%$ ajustado vs.\ $86{,}17\%$ ABG, $85{,}37\%$ \texttt{ipa-dict} e $84{,}17\%$ WikiPron). Experimentos complementares em trabalho paralelo (inglês, russo, hebraico, turco e francês) trazem evidência adicional da portabilidade da arquitetura entre três famílias linguísticas (indo-europeia, afro-asiática e túrquica) e três sistemas de escrita.
\end{resumo}

\palavraschave{correção ortográfica, português brasileiro, conversão grafema-fone, traços distintivos, distância fonética}

\begin{abstract}
Conventional spell checkers rely either on grapheme-level edit distance and keyboard proximity---failing on cognitive and phonetic errors---or on English-centric phonetic codes such as Soundex and Metaphone. In this work, we detail, adapt, and thoroughly evaluate for Brazilian Portuguese \textbf{LipSpell}, a language-parameterizable phonetic scoring framework originally introduced in companion work, combining a precomputed $|P|\times|P|$ phone-distance matrix built from universal IPA distinctive features, ABNT2 keyboard geometry with Unicode NFD (Normalization Form D) diacritic tolerance, and corpus frequency priors, accelerated by symmetric-delete and seven-class articulatory-skeleton indexing. On \textit{Wikipédia: Lista de erros comuns} ($n=1{,}149$ pairs and $n=500$ slice), LipSpell outperforms Norvig, Hunspell, and Aspell ($73.20\%$ in direct transfer and $76.20\%$ with NFD keyboard tolerance vs.\ $66.40\%$/$64.20\%$/$64.00\%$ on $n=500$; $72.24\%$ and $74.93\%$ vs.\ $66.23\%$/$63.19\%$/$60.14\%$ on $n=1{,}149$), with major gains on diacritic-only errors ($+11.62$ pp over Norvig) and multi-edit errors ($+15.34$ pp). Comparing rule-based \texttt{PortugueseG2P} against three pronunciation dictionaries (ABG, \texttt{ipa-dict}/Wiktionary, and WikiPron) within the full pipeline shows that, with grave-vowel and word-final \texttt{-z} rules, pure rules match ABG in Top-1 on the full slices ($76.20\%$ on $n=500$; $74.93\%$ on $n=1{,}149$), achieve the highest Top-5 ($85.20\%$ and $83.90\%$), and outperform all three dictionaries on clean non-word splits ($87.17\%$ tuned vs.\ $86.17\%$ ABG, $85.37\%$ \texttt{ipa-dict}, and $84.17\%$ WikiPron). Companion benchmarks in English, Russian, Hebrew, Turkish, and French provide further evidence of cross-lingual portability across three language families and three writing systems.
\end{abstract}

\keywords{spelling correction, brazilian portuguese, grapheme-to-phoneme, distinctive features, phonetic distance}

\section{Introdução}
Erros ortográficos humanos dividem-se principalmente em duas categorias \citep{kukich1992techniques,toutanova2002pronunciation}: \emph{erros tipográficos}, decorrentes de desvios motores no teclado, e \emph{erros cognitivos ou fonéticos}, nos quais o escrevente representa graficamente uma pronúncia aproximada da palavra-alvo (por exemplo, \textit{escessão} $\rightarrow$ \textit{exceção}, \textit{pesquiza} $\rightarrow$ \textit{pesquisa}, ou em inglês \textit{nite} $\rightarrow$ \textit{night}).

A partir do arcabouço \textbf{LipSpell} --- originalmente proposto e avaliado em perspectiva multilíngue em \citet{araujo2026lipspell} ---, nossa contribuição neste trabalho consiste em detalhar, adaptar e avaliar em profundidade sua aplicação ao português brasileiro, integrando distância fonética, mapa de teclado brasileiro ABNT2, distância ortográfica e frequência lexical. Enquanto \citet{araujo2026lipspell} concentra-se na formulação geral, no estudo comparativo de métricas em inglês e na transferência entre seis línguas, o presente artigo distingue-se pelo estudo linguístico e pela parametrização específica para o português brasileiro --- incluindo o conversor \texttt{PortugueseG2P} refinado com neutralização computacional da distinção de nasalidade, extensões palatais, tratamento de vogais com acento grave (\texttt{è/ì/ò/ù}) e dessonorização de \texttt{-z} em final de palavra, além da tolerância diacrítica via decomposição canônica Unicode NFD (\emph{Normalization Form D})\footnote{NFD (\emph{Normalization Form D}) é uma forma de normalização Unicode que aplica a decomposição canônica para decompor caracteres pré-compostos em seu caractere-base seguido das marcas combinantes em ordem canônica padronizada (ex.: \texttt{á} $\rightarrow$ \texttt{a}+\,\textasciiacute).} no teclado ABNT2 (Seções~\ref{sec:fonologia} e~\ref{sec:arq}) ---, além de uma avaliação experimental dedicada à correção lexical de palavras isoladas (fora do escopo de modelos contextuais de sentença para erros de palavra real). O nome \emph{LipSpell} (\textbf{L}anguage-\textbf{I}ndependent \textbf{P}honetic \textbf{Spell}ing) evoca simultaneamente os lábios (\textit{lip} em inglês, $[+\text{labial}]$) e o trocadilho fonético com \emph{leap} (ex.: \textit{lip ov fayth} $\rightarrow$ \textit{leap of faith}), em que /\textsci/ $\leftrightarrow$ /i/ e /v/ $\leftrightarrow$ /f/ diferem por apenas um traço distintivo (tensão e vozeamento).

Especificamente, investigamos três questões de pesquisa sobre o português brasileiro: (\textbf{QP1}) se a combinação de distância fonética sobre traços distintivos, geometria do teclado ABNT2 (com tolerância diacrítica NFD), distância ortográfica e frequência lexical supera corretores tradicionais (Seções~\ref{sec:pt_wiki} e~\ref{sec:comp_pt}); (\textbf{QP2}) como o conversor compacto por regras \texttt{PortugueseG2P} se comporta frente a três dicionários externos de pronúncia --- Corpus ABG, \texttt{ipa-dict} e WikiPron --- no pipeline completo (Seção~\ref{sec:pt_dicts}); e (\textbf{QP3}) em quais categorias taxonômicas de erro fonológico e ortográfico do português brasileiro concentram-se os ganhos do modelo (Seção~\ref{sec:taxonomia}).

Embora algoritmos clássicos de distância de edição \citep{damerau1964technique}, modelos de canal ruidoso \citep{church1991probability,brill2000improved} e a linha de base de \citet{norvig2007spell} tratem bem erros tipográficos, eles penalizam erros fonéticos que diferem por múltiplos grafemas. Por outro lado, \citet{toutanova2002pronunciation} modelaram pronúncias no canal ruidoso, e na fonologia computacional o algoritmo ALINE \citep{kondrak2003phonetic}, o PanPhon \citep{mortensen2016panphon} e o Epitran \citep{mortensen2018epitran} exploram traços articulatórios universais. Expandindo nosso protótipo anterior \citep{araujo2018phonetic}, integramos matrizes pré-computadas de traços distintivos, geometria de teclado e indexação tipo SymSpell \citep{garbe2012symspell}.

\section{Princípios Fonológicos e Métricas}
\label{sec:fonologia}
Na fonologia gerativa formalizada por \citet{chomsky1968sound} (\textit{The Sound Pattern of English}, SPE), os fonemas não são tratados como unidades atômicas indivisíveis, mas como feixes simultâneos de \emph{traços distintivos} subsegmentares de base articulatória (como $[+\text{voz}]$ ou $[+\text{consonantal}]$) que definem classes naturais e contrastam sons da fala por especificações binárias ($+$/$-$). Para representar o conjunto $P$ de segmentos IPA em um espaço vetorial uniforme, cada segmento $p \in P$ é codificado por uma extensão operacional ternária de 24 traços articulatórios $\mathbf{v}_p \in \{-1, 0, +1\}^{24}$ da tabela de \citet{hulden2012distinctive}, estendida com a nasal palatal /\textipa{J}/ (\texttt{nh}) e a lateral palatal /\textturny/ (\texttt{lh}) do português (Tabela~\ref{tab:ipa_features} no Anexo). Valores $+1$ e $-1$ indicam presença ou ausência do traço binário clássico, enquanto $0$ indica traço não aplicável (ex.: altura vocálica vale $0$ para consoantes); todos os 45 segmentos de $P$ possuem ao menos $15$ traços ativos ($\lVert \mathbf{v}_p \rVert_2 > 0$).

Pré-computamos uma única vez a matriz de substituição $M \in \mathbb{R}^{|P|\times|P|}$ e o vetor simétrico de inserção/remoção $\mathbf{d}_{\text{indel}} \in \mathbb{R}^{|P|}$ (acesso $O(1)$ na programação dinâmica), definido por $d_{\text{indel}}(p) = \mathrm{clip}\!\left(1{,}4\,\frac{\sum_k w_k |v_{p,k}|}{\sum_k w_k},\, 0{,}35,\, 0{,}75\right)$, de modo que vogais e aproximantes custam menos para inserir/remover do que oclusivas. Para quaisquer dois segmentos IPA $u, v \in P$ com respectivos vetores de traços $\mathbf{v}_u, \mathbf{v}_v \in \{-1, 0, +1\}^{24}$ (cujas $k$-ésimas coordenadas são $v_{u,k}$ e $v_{v,k}$), comparamos três métricas para a entrada $M_{u,v}$:
\begin{itemize}
    \item \textbf{$L_1$ uniforme\footnote{Neste artigo, $L_1$ designa a distância de Manhattan (norma $\ell_1$) normalizada entre vetores ternários de traços fonológicos --- não primeira língua (L1) nem regularização $L_1$ ---, e \emph{saliência} designa pesos fonológicos fixos $w_k$ inspirados no ALINE \citep{kondrak2003phonetic}, e não mapas de saliência de aprendizado de máquina.}:} $M_{u,v}^{L_1} = \frac{1}{2F}\sum_{k=1}^F |v_{u,k} - v_{v,k}|$ ($F=24$). Como zeros não aplicáveis contribuem com $0$, fones com muitos zeros compartilhados parecem artificialmente próximos sob $L_1$.
    \item \textbf{$L_1$ ponderado por saliência:} Inspirado no ALINE \citep{kondrak2003phonetic}, $M_{u,v}^{\text{sal}} = \frac{\sum_k w_k |v_{u,k} - v_{v,k}|}{2\sum_k w_k}$, com pesos heurísticos manuais $w_k$ por traço: $2{,}0$ (\texttt{syll}); $1{,}8$ (\texttt{cons}); $1{,}5$ (\texttt{son}, \texttt{cont}, \texttt{nas}, \texttt{lab}, \texttt{cor}, \texttt{dor}); $1{,}2$ (\texttt{approx}); $1{,}0$ (\texttt{delrel}, \texttt{labdent}, \texttt{strid}, \texttt{lat}); $0{,}8$ (\texttt{round}, \texttt{ant}, \texttt{hi}, \texttt{lo}, \texttt{fr}, \texttt{bk}); $0{,}6$ (\texttt{voi}, \texttt{dist}); e $0{,}5$ (\texttt{spread}, \texttt{constr}, \texttt{tns}).
    \item \textbf{Distância cosseno:}
    \begin{equation}
        M_{u,v}^{\cos} = \frac{1}{2}\left(1 - \frac{\mathbf{v}_u \cdot \mathbf{v}_v}{\lVert \mathbf{v}_u \rVert_2 \, \lVert \mathbf{v}_v \rVert_2}\right) \in [0, 1].
    \end{equation}
    Como $\mathbf{v}_p \in \{-1,0,+1\}^{24}$, o cosseno varia em $[-1,1]$ e o fator $\frac{1}{2}$ reescala para $[0,1]$. Na tabela de 45 segmentos, consoantes possuem em média $18{,}25 \pm 1{,}50$ traços ativos ($n=28$) contra $19{,}94 \pm 0{,}87$ das vogais ($n=17$); normalizar por $\lVert \mathbf{v}_p \rVert_2$ reduz a influência dessa diferença de magnitude, embora zeros não aplicáveis ainda possam produzir similaridades interclasses imperfeitas.
\end{itemize}
Como os custos são não uniformes e a distância cosseno não satisfaz necessariamente a desigualdade triangular, o alinhamento resultante constitui um \emph{custo mínimo de Alinhamento Ótimo de Cadeias (OSA) ponderado} (com custo fixo $0{,}45$ para transposições adjacentes restritas, sem normalização pelo comprimento da palavra), perfeitamente computável via programação dinâmica.

\section{Arquitetura Unificada e Regras G2P do Português}
\label{sec:arq}
A geometria de teclado é pré-computada em matrizes $|K|\times|K|$ para QWERTY e ABNT2 (incluindo \texttt{ç} e acentos), saturada em $3{,}0$ larguras de tecla e normalizada em $[0,1]$. Para tolerar a omissão ou troca de acentos sem penalizá-la como erro de tecla distante, aplicamos decomposição NFD, que separa cada grafema em letra-base e marca diacrítica (ex.: \texttt{á}$\rightarrow$\texttt{a}+\,\textasciiacute), atribuindo custo reduzido de $0{,}15$ quando dois caracteres compartilham a mesma letra-base (como \texttt{a}/\texttt{á} ou \texttt{e}/\texttt{ê}). Para a função de conversão grafema--fone (\textit{grapheme-to-phoneme}, G2P) $\phi$, que mapeia uma palavra escrita $w$ na sua sequência de fones IPA $\phi(w)$ (\texttt{PortugueseG2P}), adotamos uma aproximação fonêmica simplificada baseada em um conjunto determinístico de regras ordenadas --- dígrafos e ditongos, regras contextuais e mapeamentos não triviais (omitindo-se apenas as correspondências diretas $1$-para-$1$) --- organizado na Tabela~\ref{tab:pt_g2p_rules}, usando /h/ como símbolo canônico operacional para o \textit{r}-forte (como em \texttt{rr} e \texttt{r} inicial, sem pretender abranger todas as realizações regionais velares ou uvulares). Exceções lexicais e alofonias não são modeladas explicitamente: grafemas ambíguos recebem um mapeamento canônico único (\texttt{x}$\rightarrow$/\textesh/, \texttt{V\_s\_V}$\rightarrow$/z/, não distinguindo por exemplo \textit{táxi} `/ks/` ou \textit{exame} `/z/`), os dígrafos \texttt{qu} e \texttt{gu} são mapeados incondicionalmente para /k/ e /\textipa{g}/ (omitindo realizações com semivogal como \textit{quase} `/kw/` e \textit{linguiça} `/gw/`), vogais com acento grave (\texttt{è/ì/ò/ù}, não pertencentes à ortografia regular) são normalizadas para suas vogais-base para absorver ruídos de digitação como \textit{albùm}, e a nasalidade é omitida computacionalmente tanto nas vogais quanto nos ditongos nasais (\texttt{â/ã}$\rightarrow$/\textturna/, \texttt{ô/õ}$\rightarrow$/o/, \texttt{ão}$\rightarrow$/a\textupsilon/, \texttt{õe}$\rightarrow$/o\textsci/, \texttt{ãe}$\rightarrow$/a\textsci/) para evitar, na grande maioria dos casos de omissão de til ou circunflexo (\textit{nao} vs.\ \textit{não}), uma penalidade fonética artificial --- ainda que, em certos pares mínimos, essa neutralização computacional suprima um contraste fonológico real.

\begin{table*}[ht]
\centering
\scriptsize
\setlength{\tabcolsep}{2.8pt}
\caption{Regras determinísticas do conversor grafema--fone \texttt{PortugueseG2P} (ordem de prioridade: contextuais/multigrafos $\rightarrow$ simples).}
\label{tab:pt_g2p_rules}
\begin{tabular}{lp{5.1cm}ll}
\toprule
Categoria & Regra / Padrão & Saída IPA & Exemplos \\
\midrule
Dígrafos consonantais & \texttt{ch}$\rightarrow$/\textesh/, \texttt{lh}$\rightarrow$/\textturny/, \texttt{nh}$\rightarrow$/\textipa{J}/, \texttt{rr}$\rightarrow$/h/, \texttt{ss}$\rightarrow$/s/ & /\textesh, \textturny, \textipa{J}, h, s/ & \textit{chave, milho, banho, carro} \\
Dígrafos velares & \texttt{qu}$\rightarrow$/k/, \texttt{gu}$\rightarrow$/\textipa{g}/ (incondic.; omite /kw, gw/) & /k, \textipa{g}/ & \textit{queijo, guerra} \\
Ditongos (sem nasal.) & \texttt{ão}$\rightarrow$/a\textupsilon/, \texttt{õe}$\rightarrow$/o\textsci/, \texttt{ãe}$\rightarrow$/a\textsci/, \texttt{ou}$\rightarrow$/o/, \texttt{ei}$\rightarrow$/e\textsci/, \texttt{ai}$\rightarrow$/a\textsci/, \texttt{oi}$\rightarrow$/\textopeno\textsci/, \texttt{au}$\rightarrow$/a\textupsilon/ & /a\textupsilon, o\textsci, a\textsci, o, e\textsci, a\textsci, \textopeno\textsci, a\textupsilon/ & \textit{coração, põe, mãe, ouro} \\
Regras contextuais & \texttt{sc/xc}+\texttt{e,i}$\rightarrow$/s/; \texttt{c}+\texttt{e,i}$\rightarrow$/s/; \texttt{g}+\texttt{e,i}$\rightarrow$/\textyogh/; \texttt{V\_s\_V}$\rightarrow$/z/; \texttt{\#r}$\rightarrow$/h/; \texttt{z\#}$\rightarrow$/s/ & /s, \textyogh, z, h/ & \textit{nascer, cedo, gelo, casa, rato, atraz} \\
Vogais e consoantes & \texttt{á/à}$\rightarrow$/a/, \texttt{â/ã}$\rightarrow$/\textturna/, \texttt{é}$\rightarrow$/\textepsilon/, \texttt{ê/è}$\rightarrow$/e/, \texttt{í/ì}$\rightarrow$/i/, \texttt{ó}$\rightarrow$/\textopeno/, \texttt{ô/õ/ò}$\rightarrow$/o/, \texttt{ú/ù}$\rightarrow$/u/, \texttt{ç}$\rightarrow$/s/, \texttt{j}$\rightarrow$/\textyogh/, \texttt{x}$\rightarrow$/\textesh/, \texttt{h}$\rightarrow\emptyset$ & /a, \textturna, \textepsilon, e, i, \textopeno, o, u, s, \textyogh, \textesh/ & \textit{café, avô, albùm, caça, hora} \\
\bottomrule
\end{tabular}
\end{table*}

Combinamos dois índices: (1) um índice de remoções simétricas (\textit{SymDelete}/\textit{SymSpell} \citep{garbe2012symspell}), que pré-indexa todas as cadeias obtidas pela remoção de $0$ a $2$ caracteres de cada palavra do vocabulário e as cruza com as remoções de $0$ a $2$ caracteres da consulta $x$ para recuperar candidatos até distância de edição ortográfica $\le 2$, e (2) um código fonético exclusivo de sete macroclasses heurísticas de articulação (definidas operacionalmente para indexação, e não como classes naturais fonológicas estritas) --- vogais (\texttt{V}), nasais (\texttt{N}), líquidas/aproximantes (\texttt{L}), sibilantes/fricativas (\texttt{S}), labiais (\texttt{P}), coronais (\texttt{T}) e dorsais/glotais (\texttt{K}) ---, atribuídas por ordem de precedência $\texttt{V} > \texttt{N} > \texttt{L} > \texttt{S} > \texttt{P} > \texttt{T} > \texttt{K}$ (o símbolo $>$ indica prioridade na classificação, isto é, cada fone recebe apenas a primeira classe aplicável na hierarquia, como /m/ em \texttt{N} antes de \texttt{P})\footnote{Os segmentos IPA de cada macroclasse (detalhados em \citet{araujo2026lipspell}) são: $\texttt{V}$ (vogais); $\texttt{N}$ (/m, n, \textipa{J}/); $\texttt{L}$ (/l, \textturny, \textfishhookr/); $\texttt{S}$ (/s, z, \textesh, \textyogh/); $\texttt{P}$ (/p, b, f, v/); $\texttt{T}$ (/t, d/); e $\texttt{K}$ (/k, \textipa{g}, h/), funcionando de forma análoga a chaves fonéticas do tipo Soundex \citep{russell1918soundex,kukich1992techniques}, Metaphone \citep{philips1990metaphone} e Double Metaphone \citep{philips2000double}.}, colapsando classes consecutivas idênticas e indexando também as variantes com remoção de $1$ símbolo do código (para tolerar omissão ou inserção de uma classe inteira), com as listas de cada chave ordenadas por frequência decrescente (limitadas aos $300$ candidatos mais frequentes na consulta para limitar o custo de alinhamento em chaves muito povoadas).

A função de pontuação discriminativa inspirada no modelo de canal ruidoso combina custos OSA totais (não normalizados pelo comprimento): $d_{\text{orth}}(x,w)$ é a distância OSA ortográfica unitária (custo $1$ para substituição, inserção, remoção e transposição); $d_{\text{kb}}(x,w)$ usa substituição pela distância normalizada de teclado em $[0,1]$, custo de inserção/remoção dependente do contexto adjacente $\min(0{,}7,\, 0{,}4 + 0{,}6\, d_{\text{kb}})$ e transposição $0{,}45$; e $d_{\text{phon}}(\phi(x),\phi(w))$ aplica a matriz $M^{\cos}$, $\mathbf{d}_{\text{indel}}$ e transposição $0{,}45$. O termo lexical adota suavização empírica de Lidstone ($0{,}5$), $\log P(w)=\log(c(w)+0{,}5)-\log(\sum_{v\in V}c(v)+0{,}5|V|)$ sobre o vocabulário em minúsculas preservando diacríticos, com desempate determinístico por ordem lexicográfica:
{\small
\begin{equation}
\begin{aligned}
    S(w \mid x) &= \lambda_{\text{freq}} \log P(w) - \lambda_{\text{orth}} d_{\text{orth}}(x,w) \\
                &\quad - \lambda_{\text{kb}} d_{\text{kb}}(x,w) \\
                &\quad - \lambda_{\text{phon}} d_{\text{phon}}(\phi(x),\phi(w)),
\end{aligned}
\end{equation}%
}%
com pesos $\lambda_{\text{freq}}=1{,}0$, $\lambda_{\text{orth}}=4{,}0$, $\lambda_{\text{kb}}=1{,}5$ e $\lambda_{\text{phon}}=4{,}0$ ajustados previamente em uma amostra independente de $800$ pares em inglês da base Birkbeck \citep{mitton1985birkbeck} (corpus clássico de erros ortográficos compilado por Roger Mitton, sem sobreposição com nenhum conjunto de teste), sendo aplicados ao português sem ajuste específico salvo indicação explícita (divisão $50/50$). Chamamos este modelo de quatro critérios (usando distância cosseno, salvo indicação) de \textbf{LipSpell} para distingui-lo das linhas de base de critério único, enquanto o arcabouço geral de pontuação é parametrizável por idioma.

\section{Resultados Experimentais}
Todos os experimentos foram executados em Intel Core i7-13650HX ($14$\,GiB RAM, Ubuntu 26.04.1 LTS, kernel Linux 7.0, Python 3.14 + NumPy). Código e scripts estão disponíveis em \url{https://github.com/leolca/lipspell}.

\subsection{Português Brasileiro (\textit{Wikipédia: Lista de erros comuns})}
\label{sec:pt_wiki}
Indexamos $240.040$ palavras em português ($11{,}90$\,s de construção, $1.836{,}6$\,MiB de memória residente, RSS; tempo médio de consulta de $19{,}07$\,ms/palavra) combinando a lista de frequências de Conteúdo Gerado pelo Usuário (\texttt{freq-cgu.txt}) com o léxico \texttt{unitex-full-clean}, ambos distribuídos com o normalizador Enelvo \citep{bertaglia2016enelvo}\footnote{No Enelvo, os dois recursos cumprem papéis distintos e complementares: \texttt{unitex-full-clean} (derivado do léxico Unitex-PB) fornece uma lista de referência de formas válidas do português sem contagens de ocorrência, enquanto \texttt{freq-cgu.txt} traz contagens empíricas extraídas de corpus de Conteúdo Gerado pelo Usuário (CGU), com menor cobertura de formas formais e flexionadas. Combinamos ambos (filtrando entradas alfabéticas e atribuindo piso $1$ às palavras do léxico ausentes em \texttt{freq-cgu.txt}) para unir cobertura vocabular ampla e estimativas de frequência para o termo $\log P(w)$.} (atribuindo piso de frequência $1$ às entradas sem contagem explícita). Avaliamos principalmente no conjunto completo A--Z ($n=1.149$ pares únicos de \textit{Wikipédia: Lista de erros comuns} \citep{wikipedia_erros_comuns}, eliminando viés de prefixo alfabético) e, complementarmente para comparação com a análise por componentes, nos primeiros $500$ pares deduplicados em ordem alfabética. Todas as acurácias (Aspell 0.60.8.2 \citep{atkinson_aspell}, Hunspell 1.7.2 \citep{nemeth_hunspell}, Norvig \citep{norvig2007spell} e LipSpell, executados em UTF-8 preservando minúsculas e diacríticos sobre as $5$ primeiras sugestões) são calculadas sobre o denominador completo ($n=1.149$ ou $n=500$), computando palavras fora do vocabulário (OOV) como erro; os intervalos após $\pm$ indicam intervalos de confiança de $95\%$ pela aproximação normal da distribuição binomial ($1{,}96\sqrt{\hat{p}(1-\hat{p})/n}$).

Mesmo com cobertura vocabular ligeiramente inferior à do Aspell e Hunspell ($96{,}8\%$ vs.\ $98{,}4\%$ em $n=500$; $95{,}8\%$ vs.\ $97{,}4\%$ em $n=1.149$), o modelo unificado na transferência direta (sem tolerância diacrítica NFD no teclado) já atinge $73{,}20\pm3{,}88\%$ Top-1 ($85{,}20\%$ Top-5) em $n=500$ ($p \le 4{,}53\times 10^{-5}$ pelo teste exato de McNemar emparelhado; discordâncias favoráveis/desfavoráveis de $40/6$ contra Norvig, $81/35$ contra Aspell e $82/37$ contra Hunspell) e $72{,}24\pm2{,}59\%$ Top-1 em $n=1.149$ contra $66{,}23\%$ de Norvig ($82/13$, $p=2{,}08\times 10^{-13}$), $63{,}19\%$ do Hunspell ($200/96$, $p=1{,}46\times 10^{-9}$) e $60{,}14\%$ do Aspell ($219/80$, $p=4{,}21\times 10^{-16}$) (Tabela~\ref{tab:ptbr}); ao ativar a tolerância diacrítica NFD ($0{,}15$) no canal de teclado, a acurácia Top-1 sobe para $76{,}20\pm3{,}73\%$ em $n=500$ (mantendo a cobertura Top-5 de $85{,}20\%$ nessa fatia) e $74{,}93\pm2{,}51\%$ em $n=1.149$ (Tabela~\ref{tab:ptbr}). Em erros de múltiplas edições não diacríticas ($d_{\text{orth}}\ge 2$, $n=163$ em $n=1.149$), o modelo unificado atinge $42{,}33\%$ contra $26{,}99\%$ de Norvig ($+15{,}34$ pontos percentuais). Nota-se também que o modo \textit{Somente fonético} atinge apenas $39{,}20\%$ Top-1 em $n=500$ ($39{,}16\%$ em $n=1.149$), apesar de alcançar $68{,}20\%$ ($68{,}67\%$) em Top-5. Por inspeção programática dos $250$ erros Top-1 não-OOV desse modo em $n=500$, verifica-se que $213$ ($85{,}2\%$) decorrem de \emph{palavras de alta frequência} com o mesmo esqueleto fonético macroscópico ou pequena distância fonética (como \textit{acorda} para \textit{acêrca}, \textit{acessos} para \textit{acessor} ou \textit{excelente} para \textit{açoreana}): sem os critérios ortográfico e de teclado para penalizar divergências grafemáticas grandes, o ganho no termo de frequência $\log P(w)$ supera a vantagem de $d_{\text{phon}}=0$ do alvo correto e promove o competidor mais frequente ao Top-1 (permanecendo apenas os erros de palavra real como \textit{mas}/\textit{mais} fora do escopo lexical isolado).

\begin{table*}[ht]
\centering
\scriptsize
\setlength{\tabcolsep}{3.5pt}
\caption{Avaliação em português brasileiro (\textit{Wikipédia: Lista de erros comuns}, $n=500$ e conjunto completo $n=1.149$). \textit{Cobertura} indica a porcentagem das palavras-alvo presentes no léxico de cada sistema.}
\label{tab:ptbr}
\begin{tabular}{lccccc}
\toprule
Sistema & Cobertura ($500$/$1.149$) & Top-1 ($n=500$) & Top-5 ($n=500$) & Top-1 ($n=1.149$) & Top-5 ($n=1.149$) \\
\midrule
Aspell (\texttt{pt\_BR}, UTF-8)   & 98,4\% / 97,4\% & 64,00\% & 76,80\% & 60,14\% & 72,93\% \\
Hunspell (\texttt{pt\_BR}, UTF-8) & 98,4\% / 97,4\% & 64,20\% & 79,00\% & 63,19\% & 76,33\% \\
Norvig                     & 96,8\% / 95,8\% & 66,40\% & 83,00\% & 66,23\% & 81,29\% \\
Somente fonético           & 96,8\% / 95,8\% & 39,20\% & 68,20\% & 39,16\% & 68,67\% \\
\textbf{LipSpell} (transf.\ direta) & 96,8\% / 95,8\% & 73,20\% & \textbf{85,20\%} & 72,24\% & 83,72\% \\
\textbf{LipSpell (NFD)}            & 96,8\% / 95,8\% & \textbf{76,20\%} & \textbf{85,20\%} & \textbf{74,93\%} & \textbf{83,90\%} \\
\bottomrule
\end{tabular}
\end{table*}

\subsection{Dicionários de Pronúncia vs.\ G2P por Regras em Português}
\label{sec:pt_dicts}
Para avaliar o impacto de dicionários externos de pronúncia dentro do pipeline completo do LipSpell (e não como avaliação isolada de transcrição G2P), comparamos o \texttt{PortugueseG2P} contra três dicionários de pronúncia do português brasileiro (com \textit{fallback} por regras para palavras fora do dicionário): (1) \textbf{Corpus ABG} \citep{benevides2017abg} (\texttt{corpusabg\_ipa.csv}, $76.291$ formas), (2) \textbf{ipa-dict} \citep{opendictipadict} (\texttt{pt\_BR.txt}, derivado do Wiktionary, $95.933$ formas) e (3) \textbf{WikiPron} \citep{lee2020wikipron} (\texttt{por\_latn\_bz\_broad.tsv}, $188.055$ formas). Na normalização das transcrições dos dicionários, aplicamos decomposição NFD para remover marcas combinantes, acentos primário/secundário (/\textprimstress, \textsecstress/) e alongamentos (/\textlengthmark/), e remapeamos alofones fora do inventário (como /\textinvscr/$\rightarrow$/\textturnr/, /r, \textipa{;R}/$\rightarrow$/\textfishhookr/, /\textltilde/$\rightarrow$/l/, /\textgamma/$\rightarrow$/\textipa{g}/) aos segmentos da tabela de traços; palavras ausentes do dicionário recorrem diretamente ao \texttt{PortugueseG2P}. Avaliamos também uma partição limpa de erros não-palavra (\textit{non-word}) dentro do vocabulário: filtramos os $1.149$ pares mantendo apenas aqueles com forma incorreta ausente e alvo presente no léxico ($n=997$), embaralhamos com \texttt{Random(42)} e dividimos ao meio em desenvolvimento ($n=498$) e teste ($n=499$), comparando os pesos fixos do inglês (\textit{50/50 Fixo})\footnote{Adotam-se os pesos fixos calibrados no inglês (Birkbeck) como configuração padrão por dois motivos: (1) avaliar a transferência direta (\textit{zero-shot}) entre idiomas sem risco de sobreajuste às bases de teste em português ($n=500$ e $n=1.149$), aproveitando a normalização das escalas de distância; e (2) verificar empiricamente se o reajuste específico para o português traz ganho substancial --- comparação apresentada nas duas últimas colunas da Tabela~\ref{tab:pt_dicts} (\textit{50/50 Fixo} vs.\ \textit{50/50 Ajustado}), onde o reajuste traz apenas $+0{,}40$\,pp nas regras puras (de $86{,}77\%$ para $87{,}17\%$) e pequenas quedas nos dicionários ($-0{,}20$ a $-1{,}00$\,pp).} contra o reajuste exclusivo dos quatro pesos escalares $\boldsymbol{\lambda}$ na metade de desenvolvimento (\textit{50/50 Ajustado}, mantendo regras G2P e léxico inalterados). Conforme a Tabela~\ref{tab:pt_dicts}, na versão anterior do \texttt{PortugueseG2P} reportada em \citet{araujo2026lipspell} (sem vogais graves e \texttt{-z} final, com $73{,}00\%$/$75{,}80\%$ em $n=500$ e $86{,}97\%$ ajustado em $50/50$), o dicionário ABG apresentava uma pequena vantagem de $+0{,}40$\,pp em Top-1 sobre as regras em $n=500$. Isolando as duas adições na ablação das regras, o mapeamento de vogais graves (\texttt{è/ì/ò/ù}) responde integralmente pelo ganho de $+0{,}40$\,pp em $n=500$ ($75{,}80\%\rightarrow 76{,}20\%$) e $+0{,}17$\,pp em $n=1.149$ ($74{,}76\%\rightarrow 74{,}93\%$), enquanto a dessonorização de \texttt{-z} final alinha a representação fonológica de \texttt{-z} à de \texttt{-s} e eleva a divisão limpa ajustada para $87{,}17\%$. Assim, as regras puras igualam o ABG em Top-1 nas fatias completas ($76{,}20\%$ em $n=500$ e $74{,}93\%$ em $n=1.149$), mantêm o maior Top-5 ($85{,}20\%$ e $83{,}90\%$ vs.\ $84{,}40\%$ e $83{,}55\%$) e superam os três dicionários em Top-1 na divisão limpa $50/50$ ($86{,}77\%$ fixo --- empatado em Top-5 com o ABG em $96{,}19\%$ --- e $87{,}17\%$ ajustado). Repetindo a partição $50/50$ com cinco sementes aleatórias (\texttt{42, 0, 1, 2, 3}), o \texttt{PortugueseG2P} confirma estabilidade elevada: $87{,}17\pm0{,}86\%$ Top-1 ($97{,}03\pm0{,}43\%$ Top-5) com pesos fixos e $88{,}50\pm1{,}02\%$ Top-1 ($97{,}15\pm0{,}36\%$ Top-5) com pesos ajustados. Ressalva-se que essa avaliação mede a compatibilidade de ponta a ponta no corretor (evitando desencontros entre transcrições estreitas de dicionário e o \textit{fallback} por regras para a consulta errada), e não a qualidade fonética intrínseca das transcrições fora do pipeline.

\begin{table*}[ht]
\centering
\scriptsize
\setlength{\tabcolsep}{3.5pt}
\caption{G2P por regras vs.\ dicionários de pronúncia externos (ABG \citep{benevides2017abg}, \texttt{ipa-dict} \citep{opendictipadict} e WikiPron \citep{lee2020wikipron}) com \textit{fallback} por regras em português brasileiro (Top-1 / Top-5 \%).}
\label{tab:pt_dicts}
\begin{tabular}{lcccc}
\toprule
Conversor G2P & PT-Wiki ($n=500$) & PT-Wiki ($n=1.149$) & 50/50 Fixo ($n=499$) & 50/50 Ajustado ($n=499$) \\
\midrule
Somente regras (\texttt{PortugueseG2P}) & \textbf{76,20 / 85,20} & \textbf{74,93 / 83,90} & \textbf{86,77 / 96,19} & \textbf{87,17 / 96,79} \\
ABG (76k)      & \textbf{76,20} / 84,40 & \textbf{74,93} / 83,55 & 86,37 / \textbf{96,19} & 86,17 / 96,39 \\
\texttt{ipa-dict} (96k)    & 74,00 / 84,20 & 74,06 / 83,55 & 85,77 / 95,79 & 85,37 / 95,99 \\
WikiPron (188k) & 75,20 / 84,40 & 74,33 / 83,38 & 85,17 / 95,59 & 84,17 / 96,19 \\
\bottomrule
\end{tabular}
\end{table*}

\subsection{Taxonomia e Erros Mais Frequentes no Português Brasileiro}
\label{sec:taxonomia}
Analisando os $1.149$ pares da base \textit{Wikipédia: Lista de erros comuns} (Tabela~\ref{tab:pt_error_taxonomy}), classificamo-los de forma hierárquica e mutuamente exclusiva em quatro categorias: (1) \textbf{erros apenas de diacrítico} ($n=628$, $54{,}66\%$: palavras que se tornam idênticas após remoção de diacríticos via NFD --- abrangendo acentos gráficos, til e cedilha, ex.: \textit{acessivel}$\rightarrow$\textit{acessível}, \textit{acêrca}$\rightarrow$\textit{acerca}, \textit{nao}$\rightarrow$\textit{não}, \textit{cacar}$\rightarrow$\textit{caçar}); (2) \textbf{substituições de 1 grafema} ($n=222$, $19{,}32\%$: distância de edição Damerau--Levenshtein $d_{\text{orth}}=1$ com mesmo comprimento, excluindo a categoria anterior); (3) \textbf{inserções, remoções e transposições} de distância $1$ ($n=136$, $11{,}84\%$: $d_{\text{orth}}=1$ excluindo as duas categorias anteriores); e (4) \textbf{erros de múltiplas edições} ($d_{\text{orth}} \ge 2$, $n=163$, $14{,}19\%$, ex.: \textit{acessor}$\rightarrow$\textit{assessor}, \textit{aandeija}$\rightarrow$\textit{bandeja}). Os pares de confusão mais frequentes (contabilizados bidirecionalmente e de forma não mutuamente exclusiva por palavra) refletem fenômenos fonológicos e ortográficos típicos do português brasileiro: neutralização vocálica átona (\texttt{e}$\leftrightarrow$\texttt{i}: $48$ pares, ex.\ \textit{açoreana}$\rightarrow$\textit{açoriana}; \texttt{o}$\leftrightarrow$\texttt{u}: $17$ pares, ex.\ \textit{altoridade}$\rightarrow$\textit{autoridade}), neutralização de vozeamento (\texttt{s}$\leftrightarrow$\texttt{z}: $48$), homofonia contextual de sibilantes e palatais (\texttt{g}$\leftrightarrow$\texttt{j}: $37$; \texttt{ç}$\leftrightarrow$\texttt{s}: $13$; \texttt{c}$\leftrightarrow$\texttt{s}: $11$) e ambiguidade grafemática do \texttt{x} (\texttt{x}$\leftrightarrow$\texttt{s}: $20$), além de omissões/inserções de \texttt{s} ($37$), \texttt{h} mudo ($21$), \texttt{n} ($16$) e \texttt{r} ($14$). Como mostra a Tabela~\ref{tab:pt_error_taxonomy} --- na qual comparamos o LipSpell especificamente contra o Norvig por este ter obtido o melhor resultado entre as linhas de base em português ($66{,}23\%$ vs.\ $63{,}19\%$ do Hunspell e $60{,}14\%$ do Aspell em $n=1.149$, além de compartilhar exatamente o mesmo léxico de $95{,}8\%$ de cobertura) ---, o LipSpell supera o Norvig sobretudo nos erros apenas de diacrítico ($+11{,}62$ pp) e nos erros de múltiplas edições ($+15{,}34$ pp), onde a distância fonética recupera grafias foneticamente equivalentes porém ortograficamente distantes; na categoria de inserções, remoções e transposições de distância $1$, o Norvig mantém ligeira vantagem ($-0{,}73$ pp) porque a distância ortográfica unitária já isola esses casos diretamente, sem perturbações de custos fonéticos de inserção. Como a lista da Wikipédia é uma compilação colaborativa dominada por desvios diacríticos ($54{,}66\%$) em vez de uma amostra ponderada de texto espontâneo, a acurácia global não deve ser generalizada diretamente para textos reais; ainda assim, excluindo os $628$ erros apenas de diacrítico ($n=521$), o LipSpell mantém vantagem sólida de $70{,}44\%$ contra $65{,}26\%$ do Norvig ($+5{,}18$ pp).

\begin{table*}[ht]
\centering
\small
\caption{Distribuição dos erros em português (\textit{Wikipédia}, $n=1.149$) e acurácia Top-1 por categoria.}
\label{tab:pt_error_taxonomy}
\begin{tabular}{lrrccc}
\toprule
Categoria de Erro & $n$ & Proporção (\%) & Norvig (\%) & LipSpell (\%) & Ganho (pp) \\
\midrule
Apenas diacrítico          & 628 & 54,66 & 67,04 & \textbf{78,66} & $+11{,}62$ \\
Substituição ($d=1$)       & 222 & 19,32 & 81,08 & \textbf{82,43} & $+1{,}35$ \\
Ins./Rem./Transp. ($d=1$)  & 136 & 11,84 & \textbf{85,29} & 84,56 & $-0{,}73$ \\
Múltiplas edições ($d\ge 2$)& 163 & 14,19 & 26,99 & \textbf{42,33} & $+15{,}34$ \\
\bottomrule
\end{tabular}
\end{table*}

\subsection{Contribuição de Componentes, Revocação e Comparação com o Estudo Multilíngue}
\label{sec:comp_pt}
As Tabelas~\ref{tab:recall_pt} e~\ref{tab:comp_pt} isolam a geração de candidatos e a contribuição de cada critério de pontuação especificamente para o português brasileiro (\textit{PT-Wiki}, $n=500$). O código fonético de sete classes sozinho atinge $89{,}20\%$ de revocação (superando os $86{,}80\%$ das remoções simétricas \textit{SymDelete}), e a união dos dois índices alcança $96{,}40\%$ avaliando em média $280{,}5$ candidatos por consulta. No isolamento sistemático da contribuição de cada critério de pontuação/ranqueamento (Tabela~\ref{tab:comp_pt}), todas as configurações \textbf{A}--\textbf{E} compartilham exatamente a mesma amostra ($n=500$), o mesmo conjunto combinado de candidatos (mantendo constante a geração em dois estágios), os mesmos pesos e o mesmo desempate lexicográfico (enquanto o modo \textit{Somente fonético} consta na Tabela~\ref{tab:ptbr}, com $39{,}20\%$ Top-1 e $68{,}20\%$ Top-5, e a configuração \textbf{D} remove o prior de frequência para isolar o efeito conjunto dos três sinais de distância). A adição da distância fonética sobre traços distintivos (\textbf{A}$\rightarrow$\textbf{C}) produz o maior salto individual de acurácia ($+12{,}80$ pontos percentuais em Top-1, de $61{,}00\%$ para $73{,}80\%$), enquanto a combinação completa (\textbf{E}) atinge o melhor Top-5 ($85{,}20\%$ vs.\ $85{,}00\%$ de \textbf{C}). Na transferência direta sem tolerância NFD no teclado (Tabela~\ref{tab:comp_pt}), \textbf{C} fica $0{,}60$ pp acima de \textbf{E} em Top-1 ($73{,}80\%$ vs.\ $73{,}20\%$) porque vogais acentuadas recebem custo de tecla distante; com a tolerância diacrítica NFD ($0{,}15$), \textbf{B} sobe para $68{,}40\%$ e \textbf{E} atinge $76{,}20\%$ Top-1 ($85{,}20\%$ Top-5), superando todas as demais configurações em ambas as métricas.

\begin{table}[ht]
\centering
\scriptsize
\setlength{\tabcolsep}{4pt}
\caption{Revocação da geração de candidatos (\%) e média de candidatos por consulta em português (\textit{PT-Wiki}, $n=500$).}
\label{tab:recall_pt}
\begin{tabular}{lc}
\toprule
Métrica / Estágio & PT-Wiki ($n=500$) \\
\midrule
SymDelete apenas (\%)      & 86,80 \\
Código fonético apenas (\%) & 89,20 \\
Revocação combinada (\%)   & \textbf{96,40} \\
Média de candidatos / consulta & 280,5 \\
\bottomrule
\end{tabular}
\end{table}

\begin{table}[ht]
\centering
\scriptsize
\setlength{\tabcolsep}{3pt}
\caption{Análise de contribuição de componentes do modelo unificado em português (\textit{PT-Wiki}, $n=500$).}
\label{tab:comp_pt}
\begin{tabular}{lcc}
\toprule
Configuração & Top-1 (\%) & Top-5 (\%) \\
\midrule
\textbf{A}: Ortog.\ + Frequência                    & 61,00 & 81,40 \\
\textbf{B}: Ortog.\ + Teclado + Frequência          & 63,80 & 81,80 \\
\textbf{B (NFD)}: idem + NFD ($0{,}15$) & 68,40 & 83,00 \\
\textbf{C}: Ortog.\ + Fonética + Frequência         & \textbf{73,80} & 85,00 \\
\textbf{D}: Ortog.\ + Teclado + Fonética            & 71,60 & 84,20 \\
\textbf{E}: Completo (transf.\ direta)     & 73,20 & \textbf{85,20} \\
\textbf{E (NFD)}: Completo (com NFD)       & \textbf{76,20} & \textbf{85,20} \\
\bottomrule
\end{tabular}
\end{table}

Resultados detalhados da mesma arquitetura em seis conjuntos de teste em inglês\footnote{Os seis conjuntos de teste em inglês avaliados em \citet{araujo2026lipspell} são Norvig1 ($n=270$), Norvig2 ($n=400$), Aspell ($n=516$), Wikipedia-EN ($n=300$), Holbrook ($n=300$) e TOEFL-Spell ($n=1.000$).} (incluindo $1.000$ erros de aprendizes de segunda língua (L2) do TOEFL-Spell \citep{flor2019toefl}) e na transferência para outras quatro línguas de três famílias linguísticas (russo cirílico, hebraico abjad, turco e francês) --- bem como a comparação entre as métricas $L_1$, saliência e cosseno --- são apresentados no trabalho complementar \citep{araujo2026lipspell}. Em síntese, a distância cosseno alcança o melhor Top-1/Top-5 fonético e Top-5 unificado em inglês ($78{,}52\%$ em Norvig1, $77{,}30\%$ em TOEFL-Spell), e a transferência sem ajuste de pesos ou com ajuste por língua mantém ganhos consistentes sobre Norvig e Aspell entre famílias linguísticas.

\section{Conclusão}
Detalhamos, adaptamos e avaliamos em profundidade a aplicação do arcabouço \textbf{LipSpell} \citep{araujo2026lipspell} ao português brasileiro, integrando distância de edição ponderada por traços distintivos (cosseno), geometria de teclado ABNT2 (com tolerância diacrítica NFD), distância ortográfica e frequência lexical. Respondendo às três questões de pesquisa propostas: (\textbf{QP1}) na base \textit{Wikipédia: Lista de erros comuns}, a combinação dos quatro critérios supera consistentemente Norvig, Hunspell e Aspell ($73{,}20\%$ na transferência direta e $76{,}20\%$ com tolerância NFD vs.\ $66{,}40\%$/$64{,}20\%$/$64{,}00\%$ em $n=500$; $72{,}24\%$ e $74{,}93\%$ vs.\ $66{,}23\%$/$63{,}19\%$/$60{,}14\%$ em $n=1.149$), enquanto critérios isolados como o modo exclusivamente fonético sofrem com competidores de alta frequência de mesmo esqueleto fonético ($39{,}20\%$ Top-1 vs.\ $68{,}20\%$ Top-5 em $n=500$); (\textbf{QP2}) na comparação do conversor \texttt{PortugueseG2P} contra três dicionários de pronúncia (ABG, \texttt{ipa-dict}/Wiktionary e WikiPron) no pipeline completo, o refinamento das regras com vogais graves (\texttt{è/ì/ò/ù}) e dessonorização de \texttt{-z} final faz as regras determinísticas igualarem o ABG em Top-1 nas fatias completas ($76{,}20\%$ em $n=500$ e $74{,}93\%$ em $n=1.149$), alcançarem o maior Top-5 ($85{,}20\%$ e $83{,}90\%$) e superam os três dicionários em Top-1 nas divisões limpas de erros \textit{non-word} ($86{,}77\%$ fixo e $87{,}17\%$ ajustado), evitando divergências de convenção fonética com o \textit{fallback} por regras; e (\textbf{QP3}) os ganhos concentram-se sobretudo nos erros apenas de diacrítico ($+11{,}62$ pp sobre Norvig) e de múltiplas edições ($+15{,}34$ pp), mantendo vantagem de $+5{,}18$ pp ($70{,}44\%$ vs.\ $65{,}26\%$) mesmo excluindo os desvios exclusivamente diacríticos (ressaltando-se que a lista colaborativa da Wikipédia não reflete necessariamente a distribuição de frequências de textos espontâneos), associados às confusões fonológicas e ortográficas mais frequentes do português (\texttt{e}$\leftrightarrow$\texttt{i}, \texttt{s}$\leftrightarrow$\texttt{z}, \texttt{g}$\leftrightarrow$\texttt{j}, \texttt{x}$\leftrightarrow$\texttt{s}).

Resultados complementares em seis conjuntos de teste em inglês e quatro outras línguas são detalhados em \citet{araujo2026lipspell}; trabalhos futuros para o português incluem modelagem explícita de nasalidade vocálica e exceções lexicais (como os valores de \texttt{x} e integração com Epitran \citep{mortensen2018epitran}), alinhamentos $1$-para-$2$ para ditongos/dígrafos e modelos contextuais por $n$-gramas para erros de palavra real (\textit{mas}/\textit{mais}).

\bibliography{refs}

\appendix
\section{Matriz Completa de Traços Distintivos e Macroclasses}
A Tabela~\ref{tab:ipa_features} documenta integralmente os vetores ternários de $24$ traços articulatórios ($\mathbf{v}_p \in \{-1, 0, +1\}^{24}$, abreviados como $-$, $0$, $+$) e a macroclasse heurística atribuída (\texttt{V, N, L, S, P, T, K}) para todos os segmentos IPA utilizados no trabalho: os $45$ segmentos base (tabela de \citet{hulden2012distinctive} estendida com /\textipa{J}/ e /\textturny/) e, após o traço horizontal, os $8$ segmentos adicionais empregados nas línguas do estudo complementar \citep{araujo2026lipspell}.

\begin{table*}[p]
\centering
\scriptsize
\setlength{\tabcolsep}{1.8pt}
\renewcommand{\arraystretch}{0.92}
\caption{Matriz completa de traços distintivos e macroclasse heurística (\textbf{Cl.}). Colunas: \texttt{sy}=syllabic, \texttt{co}=consonantal, \texttt{so}=sonorant, \texttt{cn}=continuant, \texttt{dr}=delayed release, \texttt{ap}=approximant, \texttt{na}=nasal, \texttt{vo}=voice, \texttt{sp}=spread glottis, \texttt{ct}=constricted glottis, \texttt{lb}=labial, \texttt{rd}=round, \texttt{ld}=labiodental, \texttt{cr}=coronal, \texttt{an}=anterior, \texttt{ds}=distributed, \texttt{st}=strident, \texttt{lt}=lateral, \texttt{do}=dorsal, \texttt{hi}=high, \texttt{lo}=low, \texttt{fr}=front, \texttt{bk}=back, \texttt{tn}=tense.}
\label{tab:ipa_features}
\begin{tabular}{lc*{24}{c}}
\toprule
Seg. & Cl. & \texttt{sy} & \texttt{co} & \texttt{so} & \texttt{cn} & \texttt{dr} & \texttt{ap} & \texttt{na} & \texttt{vo} & \texttt{sp} & \texttt{ct} & \texttt{lb} & \texttt{rd} & \texttt{ld} & \texttt{cr} & \texttt{an} & \texttt{ds} & \texttt{st} & \texttt{lt} & \texttt{do} & \texttt{hi} & \texttt{lo} & \texttt{fr} & \texttt{bk} & \texttt{tn} \\
\midrule
/p/ & P & - & + & - & - & - & - & - & - & - & - & + & - & - & - & 0 & 0 & 0 & - & - & 0 & 0 & 0 & 0 & 0 \\
/t/ & T & - & + & - & - & - & - & - & - & - & - & - & - & - & + & + & - & - & - & - & 0 & 0 & 0 & 0 & 0 \\
/k/ & K & - & + & - & - & - & - & - & - & - & - & - & - & - & - & 0 & 0 & 0 & - & + & + & - & 0 & 0 & 0 \\
/b/ & P & - & + & - & - & - & - & - & + & - & - & + & - & - & - & 0 & 0 & 0 & - & - & 0 & 0 & 0 & 0 & 0 \\
/d/ & T & - & + & - & - & - & - & - & + & - & - & - & - & - & + & + & - & - & - & - & 0 & 0 & 0 & 0 & 0 \\
/g/ & K & - & + & - & - & - & - & - & + & - & - & - & - & - & - & 0 & 0 & 0 & - & + & + & - & 0 & 0 & 0 \\
/\textipa{g}/ & K & - & + & - & - & - & - & - & + & - & - & - & - & - & - & 0 & 0 & 0 & - & + & + & - & 0 & 0 & 0 \\
/t\textesh/ & S & - & + & - & - & + & - & - & - & - & - & - & - & - & + & - & + & + & - & - & 0 & 0 & 0 & 0 & 0 \\
/d\textyogh/ & S & - & + & - & - & + & - & - & + & - & - & - & - & - & + & - & + & + & - & - & 0 & 0 & 0 & 0 & 0 \\
/f/ & P & - & + & - & + & + & - & - & - & - & - & + & - & + & - & 0 & 0 & 0 & - & - & 0 & 0 & 0 & 0 & 0 \\
/\texttheta/ & T & - & + & - & + & + & - & - & - & - & - & - & - & - & + & + & + & - & - & - & 0 & 0 & 0 & 0 & 0 \\
/s/ & S & - & + & - & + & + & - & - & - & - & - & + & - & + & + & + & - & + & - & - & 0 & 0 & 0 & 0 & 0 \\
/\textesh/ & S & - & + & - & + & + & - & - & - & - & - & + & - & + & + & - & + & + & - & - & 0 & 0 & 0 & 0 & 0 \\
/h/ & P & - & - & - & + & + & - & - & - & + & - & + & - & + & - & 0 & 0 & 0 & - & - & 0 & 0 & 0 & 0 & 0 \\
/v/ & P & - & + & - & + & + & - & - & + & - & - & + & - & + & - & 0 & 0 & 0 & - & - & 0 & 0 & 0 & 0 & 0 \\
/\textipa{D}/ & T & - & + & - & + & + & - & - & + & - & - & - & - & - & + & + & + & - & - & - & 0 & 0 & 0 & 0 & 0 \\
/z/ & S & - & + & - & + & + & - & - & + & - & - & + & - & + & + & + & - & + & - & - & 0 & 0 & 0 & 0 & 0 \\
/\textyogh/ & S & - & + & - & + & + & - & - & + & - & - & + & - & + & + & - & + & + & - & - & 0 & 0 & 0 & 0 & 0 \\
/m/ & N & - & + & + & - & 0 & - & + & + & - & - & + & - & - & - & 0 & 0 & 0 & - & - & 0 & 0 & 0 & 0 & 0 \\
/n/ & N & - & + & + & - & 0 & - & + & + & - & - & - & - & - & + & + & - & - & - & - & 0 & 0 & 0 & 0 & 0 \\
/\textipa{N}/ & N & - & + & + & - & 0 & - & + & + & - & - & - & - & - & - & 0 & 0 & 0 & - & + & + & - & 0 & 0 & 0 \\
/\textturnr/ & L & - & - & + & + & 0 & + & - & + & - & - & - & - & - & + & - & + & - & - & - & 0 & 0 & 0 & 0 & 0 \\
/j/ & L & - & - & + & + & 0 & + & - & + & - & - & - & - & - & - & 0 & 0 & 0 & - & + & + & - & + & - & + \\
/l/ & L & - & + & + & + & 0 & + & - & + & - & - & - & - & - & + & + & - & - & + & - & 0 & 0 & 0 & 0 & 0 \\
/\textturnw/ & P & - & - & - & + & + & - & - & - & + & - & + & + & - & - & 0 & 0 & 0 & - & + & + & - & - & + & + \\
/w/ & L & - & - & + & + & 0 & + & - & + & - & - & + & + & - & - & 0 & 0 & 0 & - & + & + & - & - & + & + \\
/i/ & V & + & - & + & + & 0 & + & - & + & - & - & - & - & - & - & 0 & 0 & 0 & - & + & + & - & + & - & + \\
/u/ & V & + & - & + & + & 0 & + & - & + & - & - & + & + & - & - & 0 & 0 & 0 & - & + & + & - & - & + & + \\
/\textsci/ & V & + & - & + & + & 0 & + & - & + & - & - & - & - & - & - & 0 & 0 & 0 & - & + & + & - & + & - & - \\
/\textupsilon/ & V & + & - & + & + & 0 & + & - & + & - & - & + & + & - & - & 0 & 0 & 0 & - & + & + & - & - & + & - \\
/e/ & V & + & - & + & + & 0 & + & - & + & - & - & - & - & - & - & 0 & 0 & 0 & - & + & - & - & + & - & + \\
/o/ & V & + & - & + & + & 0 & + & - & + & - & - & + & + & - & - & 0 & 0 & 0 & - & + & - & - & - & + & + \\
/\textepsilon/ & V & + & - & + & + & 0 & + & - & + & - & - & - & - & - & - & 0 & 0 & 0 & - & + & - & - & + & - & - \\
/\textschwa/ & V & + & - & + & + & 0 & + & - & + & - & - & - & - & - & - & 0 & 0 & 0 & - & + & - & - & - & - & - \\
/\textturna/ & V & + & - & + & + & 0 & + & - & + & - & - & - & + & - & - & 0 & 0 & 0 & - & + & - & + & - & - & - \\
/\textrevepsilon/ & V & + & - & + & + & 0 & + & - & + & - & - & - & - & - & - & 0 & 0 & 0 & - & + & - & - & - & - & - \\
/\textturnv/ & V & + & - & + & + & 0 & + & - & + & - & - & - & - & - & - & 0 & 0 & 0 & - & + & - & - & - & + & - \\
/\textopeno/ & V & + & - & + & + & 0 & + & - & + & - & - & + & + & - & - & 0 & 0 & 0 & - & + & - & - & - & + & - \\
/\ae/ & V & + & - & + & + & 0 & + & - & + & - & - & - & - & - & - & 0 & 0 & 0 & - & + & - & + & + & - & 0 \\
/\textscripta/ & V & + & - & + & + & 0 & + & - & + & - & - & - & - & - & - & 0 & 0 & 0 & - & + & - & + & - & + & 0 \\
/\textturnscripta/ & V & + & - & + & + & 0 & + & - & + & - & - & + & + & - & - & 0 & 0 & 0 & - & + & - & + & - & + & 0 \\
/a/ & V & + & - & + & + & 0 & + & - & + & - & - & - & - & - & - & 0 & 0 & 0 & - & + & - & + & - & - & 0 \\
/\textrhookschwa/ & V & + & - & + & + & 0 & + & - & + & - & - & - & - & - & + & + & + & - & - & + & - & - & - & - & - \\
/\textipa{J}/ & N & - & + & + & - & 0 & - & + & + & - & - & - & - & - & - & 0 & 0 & - & - & + & + & - & + & - & 0 \\
/\textturny/ & L & - & + & + & + & 0 & + & - & + & - & - & - & - & - & - & 0 & 0 & - & + & + & + & - & + & - & 0 \\ \midrule
/y/ & V & + & - & + & + & 0 & + & - & + & - & - & + & + & - & - & 0 & 0 & 0 & - & + & + & - & + & - & + \\
/\o/ & V & + & - & + & + & 0 & + & - & + & - & - & + & + & - & - & 0 & 0 & 0 & - & + & - & - & + & - & + \\
/\textturnm/ & V & + & - & + & + & 0 & + & - & + & - & - & - & - & - & - & 0 & 0 & 0 & - & + & + & - & - & + & + \\
/\textfishhookr/ & L & - & + & + & + & 0 & + & - & + & - & - & - & - & - & + & + & - & - & - & - & 0 & 0 & 0 & 0 & 0 \\
/x/ & K & - & + & - & + & + & - & - & - & - & - & - & - & - & - & 0 & 0 & 0 & - & + & + & - & 0 & 0 & 0 \\
/ts/ & S & - & + & - & - & + & - & - & - & - & - & - & - & - & + & + & - & + & - & - & 0 & 0 & 0 & 0 & 0 \\
/\textglotstop/ & K & - & - & - & - & - & - & - & - & - & + & - & - & - & - & 0 & 0 & 0 & - & - & 0 & 0 & 0 & 0 & 0 \\
/\textbari/ & V & + & - & + & + & 0 & + & - & + & - & - & - & - & - & - & 0 & 0 & 0 & - & + & + & - & - & - & + \\
\bottomrule
\end{tabular}
\end{table*}

\end{document}
